\section{最大似然选择最能解释数据的参数}
\label{sec:11-Mathematical-Statistics/02-Point-Estimation/03}

矩估计把问题转化成``让样本特征等于总体特征，再反解参数''。它用汇总统计量做桥梁，不直接碰整个分布。但还有一条更直接的思路：在候选参数里，找一个让已观测数据``最不意外''的。

什么叫``最不意外''？拿到观测 \(x_1,\ldots,x_n\)。对每个候选 \(\theta\)，模型 \(P_\theta\) 给这组数据分配了一个概率（离散情形）或概率密度（连续情形）。如果一个 \(\theta\) 赋予这组数据的概率密度远高于另一个 \(\theta'\)，那在``解释已看到的东西''这件事上，\(\theta\) 比 \(\theta'\) 更强。让这个值最大的那个 \(\theta\)，就是最大似然估计。

给定观测 \(x = (x_1,\ldots,x_n)\)，似然函数把联合密度（或联合概率质量）视为参数的函数：
\[
L(\theta;x) = p_\theta(x).
\]
注意写法：\(x\) 固定，\(\theta\) 在变。这不是参数的概率分布——没有先验时，\(L(\theta;x)\) 对 \(\theta\) 积分不要求等于 1，也不需要归一化。似然只是说：如果真实参数是 \(\theta\)，观测到这组数据的概率（密度）有多大。把它当作 \(\theta\) 的函数，比较不同 \(\theta\) 下数据出现的相对可能性。

最大似然估计选取：
\[
\hat\theta_{\mathrm{MLE}} \in \argmax_{\theta\in\Theta} L(\theta;x).
\]
它回答的问题是：在候选模型族中，哪个参数让已经发生的事实看起来最不令人惊讶？

\subsection{为什么总是转到对数似然}

iid 样本的似然是 \(n\) 项乘积：
\[
L(\theta;x) = \prod_{i=1}^n p_\theta(x_i).
\]
每一项都是 0 到 1 之间的数（密度可以大于 1，但乘积在 \(n\) 大时以极高概率趋近于零）。\(0.6^{100} \approx 6.5\times 10^{-23}\)——还没算完就已经下溢成零，浮点精度完全兜不住。数值上，概率乘积是灾难。

取对数，最大化点不变——\(\log\) 严格单调递增：
\[
\ell(\theta;x) = \log L(\theta;x) = \sum_{i=1}^n \log p_\theta(x_i).
\]
乘积变求和。没有下溢。求导时：
\[
\frac{\partial}{\partial\theta}\prod_i p_\theta(x_i)
\quad\text{vs}\quad
\frac{\partial}{\partial\theta}\sum_i \log p_\theta(x_i).
\]
前者是 \(n\) 项乘积法则的噩梦，后者是 \(n\) 项简单相加。所有实际计算都在对数尺度上完成。

似然的直觉不是概率本身：在连续模型中，任何一组具体 \(x\) 的概率都严格为零（积分在零测集上消失）。似然比较的是密度高度——两个 \(\theta\) 在观测点上的相对密度之比。如果 \(L(\theta_1;x)/L(\theta_2;x) = 10\)，意思是``在这组数据上，\(\theta_1\) 的解释力是 \(\theta_2\) 的 10 倍''。这不是概率的比值，而是证据的权重。

\subsection{Bernoulli MLE}

设 \(X_i \sim \operatorname{Bernoulli}(\theta)\)，令 \(S = \sum_i X_i\)，则：
\[
L(\theta) = \theta^S (1-\theta)^{n-S}.
\]
对 \(0<\theta<1\)：
\[
\ell(\theta) = S\log\theta + (n-S)\log(1-\theta).
\]
Score（对数似然的一阶导）：
\[
\ell'(\theta) = \frac{S}{\theta} - \frac{n-S}{1-\theta}.
\]
令其为零：
\[
\frac{S}{\theta} = \frac{n-S}{1-\theta}
\;\Longrightarrow\;
\hat\theta = \frac{S}{n} = \bar X.
\]
结果与矩估计相同。在 Bernoulli 模型上，两条构造路径汇合到同一个估计量——不是因为一种方法包含了另一种，而是因为这个模型恰好让矩方程（匹配均值）和似然方程（score 零点）重叠。

二阶导数为负，内部驻点是极大值。但注意边界：如果 \(S=0\)（全不合格）或 \(S=n\)（全合格），内部驻点不存在——score 不会在 \((0,1)\) 内过零。MLE 在边界上需要直接比较端点值 \(\theta=0\) 和 \(\theta=1\) 处的似然。边界是完整优化问题的一部分，不能只依赖在开区间内求导。

\subsection{正态 MLE 与``分母为什么是 \(n\)''}

对 \(N(\mu,\sigma^2)\) iid 样本：
\[
\ell(\mu,\sigma^2) = -\frac{n}{2}\log(2\pi\sigma^2)
                     - \frac{1}{2\sigma^2}\sum_i (x_i-\mu)^2.
\]
对 \(\mu\) 求偏导：
\[
\frac{\partial\ell}{\partial\mu} = \frac{1}{\sigma^2}\sum_i (x_i-\mu),
\]
令零得 \(\hat\mu = \bar x\)。均值的 MLE 不依赖方差——在正态模型下，不管 \(\sigma^2\) 是多少，最佳猜测始终是样本均值。

代入 \(\hat\mu\) 后对 \(\sigma^2\) 求偏导（把 \(\sigma^2\) 作为一个整体参数，不分两步取平方根）：
\[
\frac{\partial\ell}{\partial\sigma^2}
= -\frac{n}{2\sigma^2} + \frac{1}{2\sigma^4}\sum_i (x_i-\bar x)^2,
\]
令零：
\[
\hat\sigma^2_{\mathrm{MLE}} = \frac{1}{n}\sum_i (x_i-\bar x)^2.
\]
分母是 \(n\)，不是 \(n-1\)。

这不是笔误，也不是``近似''。MLE 的使命是最大化似然——找到让已观测数据在模型下最不意外的参数值。无偏性不在它的目标函数里。由恒等式：
\[
\sum_i (X_i-\bar X)^2 = \sum_i (X_i-\mu)^2 - n(\bar X-\mu)^2,
\]
取期望得 \(n\sigma^2 - \sigma^2 = (n-1)\sigma^2\)，因此：
\[
\E[\hat\sigma^2_{\mathrm{MLE}}] = \frac{n-1}{n}\sigma^2,
\]
有限样本中略向下偏。使用 \(n-1\) 得无偏估计，但那个估计不再是原似然的最大点。

``MLE''和``无偏''是两条不同的评价原则，不能用一个术语替代另一个。MLE 追求似然最优——它不关心期望偏差；无偏追求期望准确——它不关心似然排序。正态方差恰好让两个目标给出不同答案，这种分歧在其他模型中也反复出现。

\subsection{Score、曲率与信息}

Score 定义为对数似然的梯度：
\[
U(\theta) = \nabla_\theta \ell(\theta).
\]
在内部极大点，\(U(\hat\theta) = 0\)——score 方程是 MLE 的一阶必要条件。

负 Hessian：
\[
-\nabla_\theta^2 \ell(\theta)
\]
衡量对数似然在 \(\theta\) 处的曲率。峰越尖，参数在局部越容易被数据区分——稍微偏离极大点，似然就迅速下降。曲率大意味着数据``锁定''参数的能力强，曲率平意味着数据对参数的区分力弱。

这个曲率（取期望后）就是 Fisher 信息，它是 MLE 渐近方差的逆——峰越尖，估计越精确。下一单元将严格展开这一串关系。

MLE 有一个值得一提的不变性：如果 \(\eta = g(\theta)\) 是一一参数变换，那么 \(\hat\eta_{\mathrm{MLE}} = g(\hat\theta_{\mathrm{MLE}})\)。因为换参数只是给同一个分布换标签，似然的排序不变。这不是数学巧合——它意味着你不需要为不同的参数化重新优化。Bayesian 的后验众数（MAP）一般不具有这个性质，因为先验密度在参数变换下不保持。

\subsection{大样本优良性质有前提}

在正则条件下——真参数在参数空间内部、模型光滑、可识别、Fisher 信息有限且非奇异——MLE 一致且渐近正态：
\[
\sqrt{n}(\hat\theta - \theta_0) \overset{d}{\longrightarrow} N(0, I(\theta_0)^{-1}),
\]
其中 \(I(\theta_0)\) 为 Fisher 信息矩阵。这一结果把 MLE 放在了渐近最优的位置：在正则条件下，不存在一致估计量能有比 \(I(\theta_0)^{-1}\) 更小的渐近方差。

但这些条件不是装饰。它们会失败：
\begin{itemize}
\tightlist
\item 边界参数：真 \(\theta_0\) 在 \(\Theta\) 的边界上，展开和正态逼近失效；
\item 混合模型、隐 Markov 模型：分量标签可交换导致不可识别，或信息矩阵奇异；
\item 参数维度随 \(n\) 增长（高维回归的 \(p \gg n\)）：经典渐近理论不适用；
\item 重尾分布：矩不存在，信息可能无限或为零，常规渐近框架站不住。
\end{itemize}
``大样本下 MLE 渐近正态''是对 MLE 性质的常见总结，但它不是独立定理——是一组前提的结论。忘掉前提就只剩口号。

\subsection{似然可能无界，MLE 可能不存在}

Gaussian 混合模型中的一个典型病态：让某个分量的方差趋零并对准单个样本点，该点的密度可以任意大，似然发散到 \(+\infty\)。此时 MLE 根本不存在——没有有限的参数值达到最大似然。优化器发散不一定是代码有 bug，可能是模型本身的数学结构没有定义良好的 MLE。

似然的多峰也是常态。复杂模型的 log-likelihood 几乎必然非凸，有多个局部极大。算法停在一个驻点，不等于找到了全局 MLE。两者之间的距离，没有任何通用定理替你弥合。

\subsection{模型错设与数值优化}

如果真实分布不在模型族中，MLE 仍可能收敛——收敛到在 KL 散度意义下最接近真实分布的伪真参数。此时普通的 Fisher 信息逆不再给出正确的渐近方差。需要 sandwich variance estimator 做稳健修正——用观测信息矩阵夹住 score 外积，校正模型错设带来的方差膨胀。

数值实现中的若干提醒：
\begin{itemize}
\tightlist
\item 始终使用 log-likelihood，不要碰概率乘积——下溢是一个会真实发生的数值灾难；
\item 对概率和协方差参数使用稳定参数化——log-尺度、logit 变换、Cholesky 分解；
\item 显式检查参数空间边界与约束——不要只依赖无约束优化器加裁切；
\item 报告优化终止状态、梯度和 Hessian 条件数——区分``算法没找到''和``MLE 不存在''；
\item 多组初值、profile likelihood、对似然面的可视化检查——非凸优化的基本功。
\end{itemize}

复杂模型的 log-likelihood 面值得你花时间去理解。看不到的面，优化器替你做了决定，但你不知道它做的是不是你想做的那个决定。
